\subsection{A parameter-interval certificate for the finite head}
\label{sec:optimized-head-interval}

We now certify degree-dependent errors on the entire flux interval
$Q\ge q_0$. The parameter is the real scalar $x=Q^{-1}\in[0,X]$, where
$X=q_0^{-1}$; $x=0$ denotes the planar limit. The frozen rows are unchanged.
For each retained pair $(r,z)$ and specified rational $\delta_{r,z}\ge0$,
the verifier proves
\begin{equation}\label{eq:interval-head-claim}
 E_{r,Q}\big|_z+\delta_{r,z}H_{r,Q}\big|_z\ge0
 \qquad\text{for every integer }Q\ge q_0.
\end{equation}
Here restriction is to the highest-weight multiplicity space of deficit
$z$; rotational invariance extends the claim to the whole spin block.
The proof encloses the full real parameter interval and does not infer
\eqref{eq:interval-head-claim} from a finite set of sampled fluxes.

We first express all finite forms in fixed polynomial coordinates. Use
the canonical basis $|A]_{x,Q}=\prod_j(\sqrt{j!}\,a_j^\dagger)^{n_j}|0\rangle$
on the sphere, with $A$ a partition of total weight $M$ into $n$ parts.
The indexed creator $x_j$ from the rational verifier remains distinct from
the scalar $x=Q^{-1}$. The planar Gram entry is
$g_A=\prod_jn_j!(j!)^{n_j}$. Define the scalar polynomials and their
positive square roots by
\begin{equation}\label{eq:interval-scalar-products}
 P_j(x)=\prod_{\ell=0}^{j-1}(1-\ell x),\quad
 T_p(x)=\prod_{\ell=0}^{p-1}(1-\ell x/2),\quad
 d_j=\sqrt{P_j},\quad e_p=\sqrt{T_p},\quad
 d_A=\prod_{j\in A}d_j.
\end{equation}
For $n\le4$ and $M\le16$, let $\mathcal P_{n,M}$ be the coordinate
space of columns indexed by these complete fixed-weight occupation
states. The invertible map to the spherical weight space is
\begin{equation}\label{eq:interval-polynomial-map}
 \Phi_{n,M,Q}:\mathcal P_{n,M}\longrightarrow
       (\Sym^nU_Q)_M,
 \qquad c\longmapsto\sum_Ac_A d_A^{-1}|A]_{x,Q}.
\end{equation}
It is a polynomial-coordinate identification, not a unitary map. Its
Gram matrix is
\begin{equation}\label{eq:interval-metric}
 G_{n,M}(x)=\Phi_{n,M,Q}^\dagger\Phi_{n,M,Q}
          =\operatorname{diag}(g_A/d_A^2).
\end{equation}
All matrices with superscript $\mathrm{pol}$ below are pullbacks of
physical quadratic forms by $\Phi$, already including this metric.

Let $\widehat B_p=\sum_{a=0}^p\partial_a\partial_{p-a}$, with
$\partial_j$ the canonical polynomial derivative. Direct substitution
of the spherical pair coefficients gives the exact, well-typed identity
\begin{equation}\label{eq:interval-pair-map}
 B_{p,Q}\Phi_{n,M,Q}
 =\frac{\kappa_p}{e_p}\Phi_{n-2,M-p,Q}\widehat B_p,
 \qquad \kappa_p^2=\frac{p!}{2^{p+1}}.
\end{equation}
The highest-weight space in polynomial coordinates is the fixed rational
kernel of $L=\sum_{j\ge1}j x_{j-1}\partial_j$: the exact raising identity
is $Q^{-1/2}J_+\Phi_{n,M,Q}=\Phi_{n,M-1,Q}L$.
If $U$ is a full rational basis of this kernel, then
\begin{equation}\label{eq:interval-common-kernel}
 \ker\big[U^\dagger H^{\mathrm{pol}}_{r,z}(x)U\big]
 =\ker\big[U^\dagger H^{\mathrm{pol}}_{r,z}(0)U\big]
 \quad (0\le x\le X).
\end{equation}
This follows by applying \eqref{eq:interval-pair-map} to every pair
annihilator; all scalar factors and Gram matrices are invertible.
The pair-sandwich structure of the normal-ordered coefficients, together
with invariance of $\ker H_r$ under rotations and spin projections,
implies that the comparison form also annihilates this same kernel exactly.

The verifier chooses a fixed linear complement using principal pivot
coordinates of the exact rational planar Hamiltonian form. Its principal
matrix is positive definite by rational elimination. A further invertible
rational upper triangular change of coordinates approximately normalizes
this Hamiltonian form to the identity. It is only a choice of basis:
Decimal arithmetic proposes its entries, but these are rounded to exact
rationals and their nonzero diagonal is checked. No bound depends on
Decimal accuracy. Write the resulting canonical coordinate frame as
$U_{\rm a}$. Positivity on its image suffices, because the exact form
annihilates the complementary kernel. A block of rank zero is identically
zero and requires no positive-pivot test.

For completeness we give the direct normal-order construction used by the
interval verifier. In a fixed channel $t$, label row columns by $0$ and
$a=(p,j)$, put $i=p+j-t$, and define the exact rational covariance
\begin{equation}\label{eq:interval-row-covariance}
 C_{ab}^{(t)}=\sum_\rho y_{\rho,a}y_{\rho,b},
 \qquad h_a(x)=\frac{i!}{e_p d_i d_j}.
\end{equation}
The sum is over the frozen rows in that channel. For $n=3,4$, define the
integer row functional $\ell_{p;l_1,\ldots,l_{n-2}}$ as the vacuum
coefficient of
$\partial_{l_1}\cdots\partial_{l_{n-2}}\widehat B_p$ on the canonical
occupation basis at weight $p+\sum_s l_s$. Integer occupation
multiplicities are included. Within this channel suppress the superscript
$t$ on $C$, and write $b=(p',j')$, $i'=p'+j'-t$. The channel's
unaveraged three-body form is
\begin{align}\label{eq:interval-direct-K3}
 &\sum_a C_{0a}\frac{h_a}{e_t}
   \big(\ell_{t;i}^{\mathsf T}\ell_{p;j}
         +\ell_{p;j}^{\mathsf T}\ell_{t;i}\big)\notag\\
 &\quad+\sum_{a,b:\,i=i'}
    \frac{C_{ab}i!}{e_p e_{p'}d_jd_{j'}}
       \ell_{p;j}^{\mathsf T}\ell_{p';j'}.
\end{align}
Its four-body form is
\begin{equation}\label{eq:interval-direct-K4}
 \sum_{a,b}C_{ab}h_ah_b
       \ell_{p;j,i'}^{\mathsf T}\ell_{p';i,j'}.
\end{equation}
Terms in \eqref{eq:interval-direct-K3} belong to weight $M=p+j$;
terms in \eqref{eq:interval-direct-K4} belong to
$M=p+j+i'=(p+j)+(p'+j')-t$. Summing channels therefore gives the full
forms $\mathsf K^{\mathrm{pol}}_{3,M}$ and
$\mathsf K^{\mathrm{pol}}_{4,M}$, supported at $M\le8$ and $M\le16$.
To verify their coefficients, the physical functional
$a_{l_1}\cdots a_{l_s}B_p$ after pullback is
\begin{equation}
 \frac{\kappa_p\sqrt{l_1!\cdots l_s!}}
 {e_p d_{l_1}\cdots d_{l_s}}\,ell_{p;l_1,\ldots,l_s}.
\end{equation}
Substitution of $\lambda=y_0/\kappa_t$ and
$\alpha_{p,j}=y_{p,j}\sqrt{i!/(\kappa_p^2j!)}$ into the exact
normal-order identities cancels the square-root factorials and gives
\eqref{eq:interval-direct-K3}--\eqref{eq:interval-direct-K4}.

The remaining targets also have finite coordinate formulas:
\begin{align}\label{eq:interval-target-forms}
 H^{\mathrm{pol}}_{n,M}(x)
 &=\sum_{p=0}^M\frac{\kappa_p^2}{e_p^2}
    \widehat B_p^{\mathsf T}G_{n-2,M-p}(x)\widehat B_p,\notag\\
 \mathsf S^{\mathrm{pol}}_{3,M}
 &=H^{\mathrm{pol}}_{3,M}G_{3,M}^{-1}H^{\mathrm{pol}}_{3,M}
       -H^{\mathrm{pol}}_{3,M},\notag\\
 \mathsf S^{\mathrm{pol}}_{4,M}
 &=\sum_{p=0}^M\frac{\kappa_p^2\kappa_{M-p}^2}
                    {e_p^2 e_{M-p}^2}\,v_p^{\mathsf T}v_p,
 \qquad v_p=\widehat B_{M-p}\widehat B_p.
\end{align}
Here $v_p$ is a row functional to the vacuum; the intermediate diagonal
maps cancel in the two successive pair annihilators. The four-body target
is $(1-c_9-\mu)\mathsf S^{\mathrm{pol}}_{4,M}$.

In polynomial coordinates the rescaled spherical lowering operator is
\begin{equation}\label{eq:interval-descendant-polynomial}
 R(x)=\sum_j(1-jx)x_{j+1}\partial_j,
 \qquad C_k(x)=R(x)^kU_{\rm a},\qquad
 n_k(x)=\prod_{l=1}^k l\,[r-(2z+l-1)x].
\end{equation}
Thus $\Phi C_k/\sqrt{n_k}$ is the normalized descendant of the chosen
highest frame. With $D_3=8,D_4=16$, exact Schur averaging gives the
highest-frame form
\begin{equation}\label{eq:interval-Schur}
 \overline K_{r,z}(x)=\frac{2+x}{r+(1-2z)x}
 \sum_{M=z}^{D_r}
 \frac{C_{M-z}(x)^{\mathsf T}\mathsf K^{\mathrm{pol}}_{r,M}(x)
 C_{M-z}(x)}{n_{M-z}(x)}.
\end{equation}
This formula applies to a nonorthonormal frame because the normalized
ladder maps preserve its Gram matrix. The finite Hermitian matrix tested
by the verifier is
\begin{equation}\label{eq:interval-final-matrix}
 F_{r,z}(x)=U_{\rm a}^{\mathsf T}
  [\mathsf T^{\mathrm{pol}}_{r,z}(x)
     +\delta_{r,z}H^{\mathrm{pol}}_{r,z}(x)]U_{\rm a}
       -\overline K_{r,z}(x),
\end{equation}
where $\mathsf T_3=\mathsf S_3$ and
$\mathsf T_4=(1-c_9-\mu)\mathsf S_4$.

We specify the interval arithmetic so that the all-$Q$ statement is
replayable. Every scalar interval has integer endpoints divided by
$2^{192}$. Products and reciprocals round lower endpoints down and upper
endpoints up using integer floor and ceiling division. Square roots use
integer square roots of the appropriately scaled endpoints, with an
increment of the upper endpoint unless its square is exact. All
denominator intervals must exclude zero and all square-root intervals
must be nonnegative; otherwise verification stops.

A Taylor interval record of order $m$ encloses the exact representation
\begin{equation}\label{eq:interval-jet-definition}
 f(x)=P_f(x)+x^{m+1}\rho_f(x),\qquad
 P_f(x)=\sum_{k=0}^m c_kx^k,\qquad 0\le x\le X.
\end{equation}
It stores intervals for every fixed coefficient $c_k$ and for all values
of $\rho_f$ on the full interval. The certificates here use $m=3$.
Addition is componentwise. If $\operatorname{tail}_{>m}$ discards
degrees at most $m$ from a polynomial, the product remainder is exactly
\begin{equation}\label{eq:interval-jet-product}
 \rho_{fg}=
 \frac{\operatorname{tail}_{>m}(P_fP_g)}{x^{m+1}}
 +P_f\rho_g+P_g\rho_f+x^{m+1}\rho_f\rho_g.
\end{equation}
The polynomial quotient in the first term is well defined also at zero.
Interval Horner evaluation of this identity encloses the product
remainder on the whole interval.

For an inverse, its Taylor coefficients satisfy
\begin{equation}\label{eq:interval-jet-inverse}
 b_0=c_0^{-1},\qquad
 b_k=-c_0^{-1}\sum_{i=1}^kc_i b_{k-i},\qquad
 \rho_{1/f}=-\frac{\rho_{fP_b}}{f},\quad
 P_b=\sum_{k=0}^mb_kx^k.
\end{equation}
The last identity holds because the polynomial part of $fP_b$ is exactly
one. For the positive square root, the corresponding exact formulas are
\begin{equation}\label{eq:interval-jet-root}
 b_0=\sqrt{c_0},\qquad
 b_k=\frac{c_k-\sum_{i=1}^{k-1}b_ib_{k-i}}{2b_0},\qquad
 \rho_{\sqrt f}=\frac{\rho_f-\rho_{P_b^2}}{\sqrt f+P_b}.
\end{equation}
The remainder identity rationalizes $\sqrt f-P_b$. These identities refer
to the exact coefficients; rounding intervals may enclose a nonzero
width around a canceled coefficient without affecting the algebraic
cancellation. Structural induction using
\eqref{eq:interval-jet-product}--\eqref{eq:interval-jet-root} proves that
the final record encloses every matrix entry throughout $[0,X]$.

After symmetric averaging of the matrix records, write the exact form as
\begin{equation}
 F_{r,z}(x)=\sum_{i=0}^m A_ix^i+x^{m+1}\mathcal E(x).
\end{equation}
All coefficient and remainder matrices are symmetric. If
$\mathcal E_{ij}(x)\in[\ell_{ij},u_{ij}]$, a rigorous uniform remainder
bound is
\begin{equation}\label{eq:interval-remainder-norm}
 \gamma=X^{m+1}\max_i\sum_j
              \max\{|\ell_{ij}|,|u_{ij}|\},\qquad
 \|x^{m+1}\mathcal E(x)\|\le\gamma.
\end{equation}
The computed $\gamma$ is rounded upward. This is the maximum absolute
row-sum bound for a symmetric matrix in the chosen finite coordinate
space; it does not assume that the coordinate frame is physically unitary.

For $s=x/X\in[0,1]$, define the Bernstein control matrices
\begin{equation}\label{eq:interval-Bernstein}
 B_k=\sum_{i=0}^k\frac{\binom ki}{\binom mi}X^i A_i,
 \qquad
 \sum_{i=0}^m A_ix^i
 =\sum_{k=0}^m\binom mk s^k(1-s)^{m-k}B_k.
\end{equation}
The scalar weights are nonnegative and sum to one. The verifier proves
$B_k-\gamma I>0$ by interval $LDL^{\mathsf T}$ elimination, requiring
each pivot interval to have strictly positive lower endpoint. Hence
every actual symmetric matrix enclosed by that interval matrix is
positive definite. If a control matrix is too conservative, repeated
arithmetic averaging performs de Casteljau subdivision into the two
half intervals. The verifier requires both halves to pass, recursively,
using the same valid $\gamma$. Bounded-depth failure returns an error
rather than an unverified positivity claim.

Equations \eqref{eq:interval-remainder-norm} and
\eqref{eq:interval-Bernstein} prove $F_{r,z}(x)\ge0$ for the whole real
interval. The common-kernel identity \eqref{eq:interval-common-kernel}
extends this to the full highest space; invertible coordinate pullback
and rotational invariance give \eqref{eq:interval-head-claim}.

The standard-library script \path{optimized_head_interval.py} takes
\texttt{--qmin}, an exact rational shift JSON through \texttt{--shifts},
and an output path through \texttt{--result}. Its saved records include
the rational frames, every Taylor coefficient and remainder interval,
all Bernstein/LDL pivot intervals, and the exact shifts for the nonzero
blocks. Decimal diagnostics are not used in any decision. A complete
run checks all $9+17=26$ relative blocks; optional single-block debugging
is explicitly distinguished in the result metadata. No numerical
eigenvalue tolerance, SVD kernel, or finite-flux extrapolation enters
this certificate.
